% Generated draft: requires mathematical and source review.
\documentclass[11pt]{article}
\usepackage[margin=25mm]{geometry}
\usepackage{fontspec}
\setmainfont{Latin Modern Roman}
\usepackage{amsmath,amssymb,mathtools,unicode-math}
\setmathfont{Latin Modern Math}
\usepackage{enumitem,needspace,xurl,hyperref}
\setlength{\parindent}{0pt}
\setlength{\parskip}{5pt}
\setlength{\emergencystretch}{3em}
\newcommand{\N}{\mathbb{N}}
\newcommand{\Z}{\mathbb{Z}}
\newcommand{\Q}{\mathbb{Q}}
\newcommand{\R}{\mathbb{R}}
\newcommand{\C}{\mathbb{C}}
\newcommand{\F}{\mathbb{F}}
\newcommand{\A}{\mathbb{A}}
\newcommand{\PP}{\mathbb{P}}
\newcommand{\E}{\mathbb{E}}
\newcommand{\eps}{\varepsilon}
\newcommand{\dd}{\,\mathrm{d}}
\begin{document}
\section*{UnsolvedMath ID: 20000992}
\textbf{Source problem code:} catalogue entry
\subsection{Definitions and mathematical statement}
\textbf{Problem 27 (Adrien Kassel).} How does chip-firing on planar metric graphs relate to analysis? The source is broad and gives no analytic target. The aid chooses real Poisson theory, planar duality and integrality of metric-divisor slopes.

Fix a finite connected plane model of a compact metric graph, refining at all source atoms, with oriented edges of positive lengths $\ell_e$. Let $B$ be its vertex-edge incidence matrix (head $+1$, tail $-1$) and $L=\operatorname{diag}(\ell_e)$. For a degree-zero atomic source $D$, the weak Poisson potential $u$, unique modulo constants, satisfies
\[
 \int_\Gamma u'\varphi'=\langle D,\varphi\rangle,
 \qquad BL^{-1}B^{\mathsf T}u=D.
\]
The edge current $j=L^{-1}B^{\mathsf T}u$ minimizes $E(q)=q^{\mathsf T}Lq$ among all real routings $Bq=D$. This is the report's minimum-energy real relaxation; the Laplacian sign convention is negative outgoing slope sum.

\textbf{Reported reciprocal-length dual certificate.} Let $C$ be the oriented face-boundary matrix with one face grounded, so $\ker B=\operatorname{im}C$. For any routing $q$, the difference from the Green current is
\[
 q-j=Cv,
\]
where the face potential solves
\[
 C^{\mathsf T}LCv=C^{\mathsf T}Lq.
\]
This is a scalar Poisson equation on the planar dual with edge lengths $\ell_e^*=1/\ell_e$, hence conductances $\ell_e$. The aid gives the exact orthogonal energy identity
\[
 E(q)=E(j)+(Cv)^{\mathsf T}L(Cv).
\]
The last term is the dual Dirichlet energy. Bridges have zero face-gradient contribution; the grounded formulation also covers the trivial face space of a tree.

For an integral degree-zero divisor $D$, choose an integral routing $q$. A principal metric divisor is the Laplacian of a continuous piecewise-linear function with integer slopes, using the declared sign convention (the opposite conventional sign gives the same principal subgroup). The aid's exact criterion is
\[
 D\text{ is principal}\quad\Longleftrightarrow\quad j\in\Z^{E}
 \quad\Longleftrightarrow\quad Cv\in\Z^{E}.
\]
It is the dual gradient, not the face potential itself, that must be integral. The fixed plane model and all lengths are part of the certificate.

\textbf{Reported atom-displacement control.} On zero-mass measures use the energy-dual norm
\[
 \|\mu\|_{H^{-1}}=\sup_{\int|\varphi'|^2\leq1}|\langle\mu,\varphi\rangle|,
\]
with functions modulo constants. The aid reports an estimate controlling movement of atoms to nearby model vertices. A standard compatible form for paired moves $D=\sum_i a_i\delta_{p_i}$, $D'=\sum_i a_i\delta_{p_i'}$ along paths of lengths at most $\delta_i$ is
\[
 \|D-D'\|_{H^{-1}}\leq\sum_i|a_i|\sqrt{\delta_i}.
\]
It controls the energy error of the corresponding real Poisson solutions. The summary does not print its exact chosen norm/constant convention, and this continuity estimate does not preserve integer slopes or divisor principality. The package supplies one rigorous analytic connection rather than resolving every meaning of the source question.
\subsection{Short English statement}
The aid relates metric chip-firing to minimum-energy Poisson currents. A planar dual with reciprocal edge lengths measures the excess energy of any routing, and integrality of its correction gradient exactly tests principality of an integral divisor. An atom-movement estimate controls the real analytic approximation, with no claim that principality survives approximation.
\subsection{Sources}
\begin{itemize}
\item Local UnsolvedMath registry, permanent ID $20000992$, \texttt{catalogue entry}; canonical \texttt{aim-combinatorics-notes.json}, \texttt{notes[116]}.
\item AIM workshop ``Generalizations of chip-firing and the critical group,'' PDF Problem 27, Adrien Kassel: \url{https://aimath.org/pastworkshops/chipfiringproblems.pdf}.
\item Research attempt 1, \texttt{partial\_result}, run status \texttt{partially\_solved}, internal rigor \texttt{proved}, low novelty confidence. Explicitly machine-generated, not peer reviewed, and to be independently verified before citation. No complete external theorem references or original literature entry supplied.
\end{itemize}
\end{document}
